\phantomsection
\addcontentsline{toc}{chapter}{Rubrics}
\fancyhead[C]{{\LARGE Rubrics}}
%\section*{Preface}
\lett{I}{n} the Christian life, the Holy Communion serves as the highest mode of worship. And due to its solemnity, it is important for the priest who celebrates the Mass to keep the rubrics and their purpose in mind. At the same time, however, the rubrics are to help the Christian offer his Sacrifice with loving care and faithfulness, not to bog him down with obligations. Therefore, Mass with its rubrics and prayers are laid out in a manner that should be accessible, lest `there be more business to find out what should be read, than to read it when it be found out.'
%It is in this mode of prayer that the Christian obediently, faithfully, and joyfully observes the Sacred Scripture's exhortation to offer the morning and evening sacrifice. And, through devotion to the Minor Hours, to `pray without ceasing'.\par
%It must be remembered, before a review of the rubrics, that the purpose of the Office is to facilitate this Sacrifice. Therefore, a simpler Office or an Office devoutly offered with imprecise rubrical distinctions, is immensely more pleasing than a despondent spirit.\par
%\emph{
%}
\section*{General Rubrics}
\begin{description}
\item[Advent Ferial Commemoration] During Advent, the Feria is commemorated with the Prayers of the First Sunday of Advent.
	
\item[Amen] When the word \emph{Amen} is in italics, it should be said by the Congregation in response to the Minister's prayer.
%This modifies the 1928 rubric, in line with the 1549 and 1662, especially with consideration of Orthodox doctrine.
%Not relevant for a Missal:
%\item[Antecommunion] It may occur that a community lacks a Priest on Sunday or another day where the Service of Holy Communion is desired. In such a case, Mattins and the Great Litany should first be said. Then, the Holy Mass may be said from the \emph{Collect for Purity} until the Sermon, inclusive, concluding with the Additional Collects (p. \pageref{AdditionalCollects}).\par
%\textsc{Note,} Mattins, the Great Litany, and Antecommunion should all be led by the Minister in the same place as the Office, not at the Altar.
%From the Monastic Diurnal

\item[Commemorations] When multiple Feast Days overlap on the same day, the Office will be of the Feast Day of highest rank. The other Feast Day(s) will then be commemorated (or, in some cases, suppressed or transferred). Commemoration consists of praying the highest-ranking Feast Day's Collect, Secret, and Postcommunion first, followed by the lower ranking Feast Day(s)' Collect, Secret, and Postcommunion.\par
\textsc{Note,} Once three Collects (Secrets, Postcommunions) have been said, further Commemorations cease and are omitted that year.

\item[Conclusion of Collects] When the Collect is addressed to God the Father, the ending is `Through Jesus Christ thy Son our Lord, who liveth and reigneth with thee, in the unity of the Holy Ghost, God, throughout all ages, world without end. Amen.'\par
If in the beginning of the Collect the Son be mentioned: `Through the same Jesus Christ thy Son our Lord; who liveth and reigneth with thee in the unity of the Holy Ghost, God, throughout all ages, world without end. Amen.'\par
If at the end of the Collect the Son be mentioned: `Who liveth and reigneth with thee in the unity of the Holy Ghost, God, throughout all ages, world without end. Amen.'\par
If the Collect be addressed to the Son: `Who livest and reignest with the Father, in the unity of the Holy Ghost, God, throughout all ages, world without end. Amen.'\par
If the Collect be addressed to the Holy Spirit: `Who livest and reignest with the Father and the Son, God, throughout all ages, world without end. Amen.'\par
When the Holy Spirit is mentioned in the Collect: `In the unity of the same Holy Ghost, etc.'

\textsc{Note,} When several proper Collects, Secrets, and Postcommunions are said, only the first and the last receive the full ending.\par
	\textsc{Note also,} The \emph{Lord be with you} is only said to introduce the first and second collects.

\item[Ferial Days] The Ferial Office, that is, the Simple Office of the occurrent Season, is always said on weekdays in Lent, on Ember Days, and on Rogation Monday, unless a I or II Double occur, in which case the Office is of the Feast with Commemoration of the Feria.\par
On weekdays in Advent and between Septuagesima and Ash Wednesday, and on Common Vigils, the Ferial Office is always said unless a Double Feast or Octave occur, and then Commemoration is made of such Ferias.\par
If a Simple Octave Day occur on one of these Ferias, it is only commemorated.\par
In like manner throughout the year the Office is of the Feria on those weekdays on which there does not occur a Double Feast, an Octave, or the Office of St. Mary on Saturday.


\item[Gloria Patri] During Passiontide, the \emph{Gloria Patri} is omitted in the Aspersion of the People, Introit, and Lavabo. During the Sacred Triduum, the \emph{Gloria Patri} is wholly omitted.

\item[Latin] In much of this Book, Latin is provided alongside the English. The Latin may always be used instead of the English, unless otherwise determined by the Ordinary. When Latin is not provided, but great devotion exists for the Latin language, the Latin which serves as the base text (such as the Liber Precum Publicarum, Missale, \& Rituale) may be used, with great concern for the catechesis and participation of the Congregation.

\item[Lenten Ferial Commemoration] During Lent, from Ash Wednesday until Palm Sunday, exclusive, the Feria is commemorated with the proper Prayers of the Lenten Feria and of Ash Wednesday.
	
\item[Liturgical Seasons] For the purpose of the rubrics, these festive seasons are referenced:
    \begin{itemize}
    	\item Advent: From I Evensong of the First Sunday of Advent until I Evensong of Christmas Day, exclusive.
        \item Christmastide: From I Evensong of Christmas Day until I Evensong of Epiphany Day, exclusive.%\footnote{This is not to be confused with the `Christmas season', which can be said to end on Candlemas.}
        \item Epiphanytide: From I Evensong of Epiphany Day until I Evensong of Septuagesima Sunday, exclusive.
        \item Septuagesimatide (Pre-Lent): From I Evensong of Septuagesima Sunday until Mattins of Ash Wednesday, exclusive.
        \item Lent: From Mattins of Ash Wednesday until I Evensong of Passion Sunday, exclusive.
        \item Passiontide: From I Evensong of Passion Sunday until Mattins of Maundy Thursday, exclusive.
        \item Sacred Triduum: From Mattins of Maundy Thursday until I Evensong of Easter Sunday, exclusive.
        \item Eastertide (Paschaltide): From I Evensong of Easter Sunday until I Evensong of Trinity Sunday, exclusive.
        \item Trinitytide: From I Evensong of Trinity Sunday until I Evensong of the First Sunday of Advent, exclusive.
    \end{itemize}

\item[Octaves] The Office of an Octave is said (or Commemoration thereof made, when it is hindered by a Feast or a Sunday) through eight continuous days in Octaves of Feasts of the I Class. Octaves of Feasts of the II Class, which are Simple Octaves, are kept only on the Octave Day itself, with the rite of a Simple, unless hindered by a more worthy Office; but no notice is taken of the days within the Octave. If any Octave is not ended before Ash Wednesday, the Vigil of Pentecost, or 17 December, no notice is taken of it thenceforth.\par
A Simple Octave Day occurring within an Octave is only commemorated.


\item[Octave Days] On the Octave Day, the Office is said as on the day of the Feast, unless otherwise noted. On a Simple Octave Day the Office is Simple.
\par
On Sundays within Privileged Octaves, the Office is said as directed in the Proper of Season. On Sundays within Common Octaves, the Office is of the Sunday with Commemoration of the Octave.

\item[Office of Our Lady on Saturday] On all Saturdays---except in Advent, Septuagesimatide, Lent, and on Ember Days---unless the Office be of a Double Feast (even transferred), or of an occurrent Octave or Vigil, or of a Sunday anticipated according to the rubrics, the Office is of St. Mary.

\item[Sunday Collects] If the Office be of a Feria without a proper or indicated Collect, then the Collect of the preceding Sunday is read. Otherwise, the Collect for the Feast Day is read instead.


%\item[Plural \& Gendered Nouns and Pronouns] Sometimes there are words in the liturgy which change depending on number and gender, such as if the person being prayed for is a woman or if there are multiple people being prayed for. The Book, in accordance with English grammar, assumes the masculine singular, and if the prayer should be changed for either number or gender, the relevant words and letters are italicised.\par
%In cases where the change could only be due to number, such as prayers for Clerics (Subdeacons, Deacons, Priests, Bishops, etc.) which are---by divine law---always masculine, the relevant words are still italicised.
\end{description}


\section*{Mass Rubrics}
\begin{description}
    %\item[Commemorations] When Feast Days overlap, the highest ranking Feast Day is celebrated with commemoration of the lower ranking Feast Days, consisting of their Collect, Secret, and Postcommunion.
%    \item[Conventual Masses] %The English Missal envisions that the Mass may be said in the context of a monastery, praying the full Monastic Breviary and having multiple Masses every day. Therefore, the conventual Mass is the higher Mass of the day, with large attendance, though the priests may celebrate a second Mass for a lower ranking Feast Day of that day, according to the rubrics.\par
    %Given that the Prayer Book Tradition does not envision many Masses being said in monastic fashion, rubrics from the English Missal regarding conventual and non-conventual Masses are largely irrelevant. For the ordinary parish church, it is safe to assume the day's Mass is a non-conventual Mass.
%    Since this Book is oriented towards use in a parish church, it assumes the Mass being said is a `non-conventual Mass'.
    \item[Gloria in excelsis] The \emph{Gloria in excelsis} is said in the Mass on Feast Days of Double rank and higher. It is also said on all Sundays, except in Advent, Septuagesimatide, and Lent.
    \item[Introduction of the Epistle] The Epistle is introduced according to manner specified in the rubric, preferably using the title of the book as given in the 1611 King James Bible.
    \item[Introit] The Introit for each Mass is provided. The Introit is said as written, followed by the \emph{Gloria Patri} (except in Passiontide) and the Introit repeated up to the `\emph{Ps}' or `\emph{Cant.}'
%    \item[Manual Acts] For the manual acts of clerics during the Mass, consultation and familiarity should be made with \emph{Ritual Notes} and the English Missal, especially during Christmastide and Holy Week.
    %Rubric from English Missal, reworded:
    \item[Nicene Creed] The Nicene Creed is said after the Gospel on all Sundays throughout the year, even if the Mass is of a Feast on which it would not otherwise be said. On all Feast Days of the Apostles, St. Joseph, Our Lady, the Holy Cross, and our Lord, and through their Octaves. From Christmas until the Octave of St. John, inclusive. On the Feast of the Epiphany and its Octave. Maundy Thursday. On Feasts of Angels, St. Mary Magdalene, St. Gregory, St. Ambrose, St. Augustine of Hippo, St. Jerome, St. Hilary, St. Isidore, Pope St. Leo I, St. Bede the Venerable, St. Peter Chrysologus, St. Athanasius, St. Basil, St. Cyril of Alexandria, St. Cyril of Jerusalem, St. Ephram Syrus, St. Gregory Nazianzen, St. John Chrysostom, St. John Damascene, St. Lawrence, St. John Baptist and his Octave Day. On All Hallows' Day and through its Octave, on the Dedication of St. Saviour, and of the Holy Apostles Peter and Paul. As well as on the anniversary of the dedication of a Church, and through its Octave; on the day of consecration of a Church, or of an Altar; on the Feasts of the Patron Saint of a Church, and on the day of a Saint where his Body or a considerable Relic is preserved; on the day of the creation and coronation of the Chief Bishop, and on the anniversary of such day; on the day, and the anniversary, of the election and consecration of a bishop.\par
    Also, the Creed is said in Votive Masses, which are celebrated solemnly for a grave matter or for a public Church cause, even if they are said in violet vestments on a Sunday.
    \item[Prefaces] The Prefaces with Solemn Chant are to be used in all Masses of at least Semidouble rank, and in Votive Masses for a grave and public cause.\par
    \textsc{Note,} When the Mass is of a Feast Day and has a proper Preface, its Preface is said in preference to the Season's.\par
    \textsc{Note,} When a Mass does not have a proper Preface, and a Feast Day with a proper Preface is commemorated, the Commemoration's Preface may be said, the Feasts of Our Lord always receiving priority.\par
    \textsc{Note,} When no proper Preface is said, the Priest proceeds immediately from \emph{Almighty, Everlasting God.} to \emph{Therefore, with Angels and Archangels}.

%    \item[Reciting the Propers] When a clerk of appropriate rank is not available, a clerk of higher rank shall fulfil his role, such as the Priest chanting the Gospel or a Deacon chanting the Epistle.
    %Translation of the following rubric is from the English Missal:
    \item[Requiem Mass] On the first day of each month (outside Advent, Lent, and Eastertide) not hindered by a Semidouble or higher Feast Day, the Mass is said generally for departed Priests, Benefactors, and others. But if no day occur in the week on which it can be said, a general commemoration of the departed shall be said, except during Lent and Eastertide.\par
    \textsc{Note,} The \emph{Dies irae} is said on All Souls' Day and on the day of the burial in all sung Masses. However, in other Masses, it may either be recited or omitted at the pleasure of the celebrant.
    \item[Seasonal Prayers] Throughout the year, as indicated in the rubrics, additional Collects, Secrets, and Postcommunions are to be said, assigned for the different seasons even within Octaves or on Vigils. On Marian Feast Days (or Feast Days which include her commemoration), when the second seasonal collect would be of St. Mary, it is instead of the Holy Ghost.\par
    \textsc{Note,} The propers chosen for the Collect should match for the Secret and Postcommunion.
    \item[Signs of the Cross] In the Mass, when {\ding{64}} appears, unless otherwise indicated or manifestly evident, it indicates the manual act of making the sign of the Cross.\par
    In the Canon, if there be a cross ({\ding{64}}) in reference to Bread or Body, it indicates a sign of the cross made over the Host. Likewise, if there be a cross ({\ding{64}}) in reference to Wine or Cup, it indicates a sign of the cross made over the Cup. Otherwise, it indicates a sign of the cross made over the Gifts generally, unless otherwise stated.\par
    When the rubrics tell the Priest to make the sign of the Cross on himself, unless otherwise indicated or manifestly evident, it indicates making the sign of the Cross from forehead to breast.
    %In Advent, however, if the office of our Lady be not said on Saturday, nevertheless the principal Mass is of her, with commemoration of Advent, unless it be an Ember Day or Vigil, as above.
\end{description}